\documentclass[12pt]{article}

% Packages for math, formatting, and layout
\usepackage{amsmath}
\usepackage{amssymb}
\usepackage{geometry}
\usepackage{hyperref}
\usepackage{enumitem}

% Set page margins
\geometry{a4paper, margin=1in}

\title{\textbf{SARSA and Q-Learning Updates}}
\author{}
\date{}

\begin{document}

\maketitle

We know Temporal Update goes like this:
$$ \text{New Estimate} \leftarrow \text{Old Estimate} + \alpha [\text{Target} - \text{Old Estimate}] $$

\noindent For action-value functions:
$$ Q(S_t, A_t) \leftarrow Q(S_t, A_t) + \alpha [\text{TD Target} - Q(S_t, A_t)] $$

\subsection*{SARSA}
\begin{itemize}[label=$\bullet$]
    \item It stands for State Action Reward State Action.
    \item It updates the action-value function ($Q(S,A)$).
    \item It evaluates how good it is for an agent to take a specific action in a specific state, assuming the agent continues to follow its current policy.
\end{itemize}

\noindent\textbf{Derivation:} \\
SARSA is derived from the Bellman Expectation Equation:
$$ q_\pi(s,a) = \mathbb{E}_\pi \left[ R_{t+1} + \gamma q_\pi(S_{t+1}, A_{t+1}) \mid S_t=s, A_t=a \right] $$

\begin{itemize}[label=$\bullet$]
    \item SARSA samples the environment.
    \item Agent takes an action, observes $R_{t+1}$, lands in $S_{t+1}$ and then uses its current policy to select $A_{t+1}$.
\end{itemize}

$$ \text{TD Target}_{\text{SARSA}} = R_{t+1} + \gamma Q(S_{t+1}, A_{t+1}) $$
$$ \therefore Q(S_t, A_t) \leftarrow Q(S_t, A_t) + \alpha \left[ R_{t+1} + \gamma Q(S_{t+1}, A_{t+1}) - Q(S_t, A_t) \right] $$

\vspace{0.5cm}

\subsection*{Q-Learning}
\begin{itemize}[label=$\bullet$]
    \item Its goal is to directly approximate the optimal policy $q^*$ instead of evaluating the current policy.
\end{itemize}

\noindent Q-learning is also derived from Bellman Expectation Equation:
$$ q^*(s,a) = \mathbb{E} \left[ R_{t+1} + \gamma \max_{a'} q_\pi(S_{t+1}, a') \mid S_t=s, A_t=a \right] $$

\begin{itemize}[label=$\bullet$]
    \item Again we sample the environment.
    \item However, instead of waiting to see what action the agent actually takes next, we look at all possible actions in $S_{t+1}$ and select one with the highest estimated Q-value.
\end{itemize}

$$ \text{TD Target}_{\text{Q-Learning}} = R_{t+1} + \gamma \max_a Q(S_{t+1}, a) $$
$$ \therefore Q(S_t, A_t) \leftarrow Q(S_t, A_t) + \alpha \left[ R_{t+1} + \gamma \max_a Q(S_{t+1}, a) - Q(S_t, A_t) \right] $$



\section*{Comparison: SARSA vs. Q-Learning}

\begin{table}[h!]
\centering
\renewcommand{\arraystretch}{1.8}
\begin{tabular}{| p{0.18\textwidth} | p{0.35\textwidth} | p{0.35\textwidth} |}
\hline
\textbf{Feature} & \textbf{SARSA} & \textbf{Q-Learning} \\ \hline
\textbf{Policy Type} &  Learns the value of the policy it is currently executing. & Learns the value of the optimal policy independently of the agent's actions. \\ \hline
\textbf{TD Target} & $R_{t+1} + \gamma Q(S_{t+1}, A_{t+1})$ & $R_{t+1} + \gamma \max_a Q(S_{t+1}, a)$ \\ \hline
\textbf{Action Selection} & Requires knowing the \textit{actual} next action ($A_{t+1}$) before updating. & Only requires knowing the next state ($S_{t+1}$) to find the max action. \\ \hline
\textbf{Risk} & Learns safer policies if exploration is high. &Learns the optimal path even if it's dangerous during exploration. \\ \hline
\textbf{Performance} & Accumulates more reward during training because it avoids fatal failures. & Accumulates less reward during training if the optimal path is near severe penalties. \\ \hline
\textbf{Convergence} & Converges to the optimal policy \textit{only if} the exploration rate ($\epsilon$) slowly decays to 0. & Converges to the optimal policy as long as all state-action pairs are visited infinitely often. \\ \hline
\end{tabular}
\end{table}

\end{document}